% ------------------------------------------------------------------------------
% Capítulo 3 · Plan de gestión del proyecto
% Fuentes: docs/tfg/{PLAN,PROGRESO,BITACORA,DECISIONES}.md, AGENTS.md,
% .claude/, .github/workflows/ci.yml, git log (54 commits, 29/09–20/10/2026)
% y el borrador memoria/borradores/02-planificacion.tex.
% Fecha de corte de los datos: 20/10/2026.
% ------------------------------------------------------------------------------
\label{cap:plan}

Este capítulo recoge la gestión del proyecto: la metodología seguida, las
tecnologías elegidas, la planificación y su comparación con lo realmente
ejecutado, la organización, los costes, los riesgos, la gestión de la
configuración y el aseguramiento de la calidad.

% ---------------------------------------------------------------------
\section{Metodología de desarrollo}
\label{sec:plan-metodologia}

\subsection{Desarrollo incremental por fases}
\label{sec:plan-incremental}

El proyecto sigue un modelo de desarrollo \textbf{iterativo e
incremental}~\cite{larman2003iterative}. El trabajo se divide en nueve fases
(F0--F8) y, dentro de las fases grandes, en incrementos más pequeños (por
ejemplo, F2.1--F2.9 en la integración del CMS). Cada fase o incremento tiene,
fijados de antemano en el plan del proyecto:

\begin{itemize}
  \item un \textbf{objetivo} y una lista de entregables;
  \item un \textbf{criterio de «hecho»} verificable (por ejemplo, «la
        aplicación arranca en local, CU-01 funciona a mano y la integración
        continua está en verde»); y
  \item sus \textbf{dependencias} con otras fases.
\end{itemize}

Un incremento solo se cierra cuando supera su criterio de «hecho» y la cadena
de verificación común descrita en la sección~\ref{sec:plan-calidad}. Así, al
final de cada incremento el sistema queda en un estado funcional y comprobado,
y un retraso afecta al alcance de los incrementos siguientes, no a la
estabilidad de lo ya construido. Esta forma de trabajar encaja con el objetivo
de una arquitectura reutilizable: cada incremento añade un módulo que debe
integrarse sin romper los anteriores (RNF-04, RNF-06).

No se ha adoptado un marco ágil formal como Scrum: con un único desarrollador
no tienen sentido sus roles ni sus ceremonias. Sí se han tomado de los métodos
ágiles la entrega de incrementos funcionales cortos, la definición explícita
de «hecho», la integración continua y la revisión frecuente con el cliente,
papel que aquí desempeña el propio autor como responsable del producto: varias
decisiones de alcance (la integración del CMS en la consola, la consola por
áreas o la interfaz solo en español) surgieron de su revisión de un incremento
ya terminado.

Las decisiones de arquitectura o de alcance se registran en el momento en que
se toman como \textit{Architecture Decision Records}~\cite{nygard2011adr},
con su contexto, las alternativas consideradas y sus consecuencias. A la fecha
de corte hay veintiuno (ADR-001 a ADR-021); los relevantes para el diseño se
resumen en el capítulo de Diseño del Sistema. Las incidencias ---fallos
de código, problemas de entorno y de herramientas--- se anotan
también en el momento, en una bitácora con el formato \emph{qué pasó, causa,
solución, evidencia y lección}. La bitácora es la fuente de los riesgos
materializados (sección~\ref{sec:plan-riesgos}) y de las lecciones aprendidas
del capítulo~\ref{cap:conclusiones}.

\subsection{Desarrollo asistido por agentes de IA}
\label{sec:plan-ia}

% PENDIENTE (autor): comprobar si la ESI/UCA exige una declaración formal del
% uso de herramientas de IA generativa en el TFG (formato y ubicación).

Una parte sustancial del código, de las pruebas y de la documentación se ha
producido con la ayuda de un asistente de programación basado en modelos de
lenguaje que actúa como \textit{agente}: un programa que, a partir de una
instrucción, lee ficheros, ejecuta órdenes y propone o aplica cambios en el
repositorio. Se ha usado Claude Code~\cite{claudecode}. Los
\textit{commits} del repositorio registran esa coautoría en su mensaje.
\todo{Confirmar el modelo o modelos de IA empleados y el periodo de uso para
la declaración de uso de IA, si la normativa de la ESI la exige.}

Para que este modo de trabajo sea controlable y evaluable, antes de escribir
código de la aplicación (fase F0) se construyó un \textit{harness} o arnés de
directrices y controles automáticos que acota lo que hacen los agentes y
comprueba el resultado. Es una aportación propia del trabajo y tiene cinco
componentes:

\begin{itemize}
  \item \textbf{Directrices escritas.} Un fichero raíz de directrices para
        agentes y desarrolladores y otros por paquete, que fijan la pila, las
        convenciones, las reglas de seguridad de la base de datos, el idioma
        de los comentarios y de la documentación y una lista de verificación
        obligatoria al terminar cualquier cambio.
  \item \textbf{Documentación de seguimiento}: plan, requisitos, decisiones,
        trazabilidad, progreso con la dedicación, y bitácora. Un agente debe
        leer el progreso y el plan de la fase antes de empezar, y no puede
        reabrir una decisión sin registrar un ADR nuevo.
  \item \textbf{Procedimientos reutilizables} (\textit{skills}) para tareas
        repetitivas o delicadas: incorporar un módulo,
        reescribir sus comentarios, registrar una decisión, revisar un cambio
        de forma adversarial, auditar las políticas de seguridad de la base de
        datos o redactar un borrador de la memoria.
  \item \textbf{Controles automáticos} (\textit{hooks}): al inicio de cada
        sesión se inyectan la fase actual y sus tareas pendientes; tras cada
        edición se ejecutan sobre el fichero los controles automáticos del
        repositorio; y antes de terminar se recuerda actualizar el
        progreso si hay cambios sin registrar.
  \item \textbf{Permisos explícitos}: los agentes no pueden escribir fuera del
        repositorio, leer ficheros de secretos ni
        hacer \textit{commit} o \textit{push} sin orden expresa del autor.
\end{itemize}

\paragraph{Reparto de responsabilidades.}
\textbf{El autor} fija el alcance y los requisitos, toma las decisiones de
arquitectura ---la integración del CMS en la consola (ADR-011) o la
organización por áreas (ADR-020) fueron propuestas suyas---, revisa y acepta o
rechaza los cambios, hace pruebas manuales exploratorias y es el único que
autoriza los \textit{commits}. \textbf{Los agentes} analizan el código
existente, proponen alternativas, implementan los incrementos, escriben las
pruebas, ejecutan la verificación, redactan la documentación de seguimiento y
los borradores de la memoria, y hacen las revisiones adversariales.
\todo{El autor completa este reparto con su propia valoración: qué partes ha
escrito o reescrito a mano, qué criterio ha seguido para aceptar cambios y
cuánto tiempo ha dedicado a la revisión.}

La revisión humana sigue siendo imprescindible. La incidencia B-14 (la consola
se cargaba sin segundo factor) la detectó el autor con una prueba manual que
ninguna prueba automática cubría, y la B-11 muestra que un test en verde no
garantiza la corrección: un test existente afirmaba precisamente el
comportamiento inseguro.

\paragraph{Modo de ahorro.}
La verificación exhaustiva con varios agentes en paralelo encontró brechas
reales, pero su coste en uso de la herramienta resultó excesivo (B-48). Desde
la F2.7c se trabaja en un \emph{modo de ahorro}: un único agente por paso,
pruebas E2E acotadas a la zona tocada y la suite completa solo al cerrar una
fase, antes de fusionar. El nivel de revisión se decide por riesgo y no por
defecto; las incidencias B-51 y B-54 muestran el precio de ese ajuste y
justifican ejecutar siempre la suite completa antes de abrir una
\textit{pull request}.

% ---------------------------------------------------------------------
\section{Tecnologías}
\label{sec:plan-tecnologias}

La pila tecnológica se elige por usar un único lenguaje tipado (TypeScript) en
el cliente, el servidor y el acceso a la base de datos, por apoyarse en
Supabase para los datos, la autenticación y el almacenamiento con la
autorización aplicada en la propia base de datos, y por emplear una sola pila
de librerías en toda la web (ADR-013). La tabla~\ref{tab:plan-tecnologias} resume cada
tecnología y su justificación.

\begin{table}[htbp]
  \centering
  \footnotesize
  \caption{Tecnologías empleadas y justificación.}
  \label{tab:plan-tecnologias}
  \begin{tabularx}{\textwidth}{@{}>{\raggedright\arraybackslash}p{3.3cm}>{\raggedright\arraybackslash}p{3.1cm}X@{}}
    \toprule
    \textbf{Tecnología} & \textbf{Uso} & \textbf{Justificación} \\
    \midrule
    TypeScript~\cite{typescript} (modo estricto) & Lenguaje de todo el código & Tipos compartidos entre cliente, servidor y base de datos; errores detectados en compilación (RNF-05). \\
    React~\cite{react} y TanStack Start~\cite{tanstack-start} & Aplicación web (\textit{frontend} y \textit{server functions}) & Renderizado en servidor, rutas por fichero con TanStack Router~\cite{tanstack-router} y funciones de servidor tipadas y validadas (RF-01, RF-02). \\
    TanStack Query~\cite{tanstack-query} & Caché de datos en el cliente & Carga en servidor y caché de consultas (RNF-09). \\
    Supabase~\cite{supabase-docs} (PostgreSQL, Auth, Storage) & Base de datos, autenticación y archivos & PostgreSQL estándar con RLS como frontera de seguridad~\cite{supabase-rls}; MFA con niveles de garantía~\cite{supabase-mfa}; entorno local completo. \\
    Hono~\cite{hono} y Drizzle ORM~\cite{drizzle} & API del CMS en \texttt{/api/cms} & API ligera sobre estándares web, montable dentro del servidor de la aplicación; consultas tipadas ejecutadas con la identidad del usuario. \\
    Stripe~\cite{stripe-docs} & Suscripciones y pagos & Pasarela con modo de pruebas completo, \textit{checkout}, portal del cliente y \textit{webhooks} (RF-07). \\
    Tailwind CSS~\cite{tailwind} y shadcn/ui~\cite{shadcn} & Estilos y componentes & Componentes accesibles, copiados al repositorio y adaptables; tema claro y oscuro (RNF-08). \\
    TanStack Form y Zod & Formularios y validación & Una sola forma de construir formularios en toda la web, incluido el CMS (ADR-013). \\
    use-intl & Internacionalización & Todos los textos en catálogos de mensajes; base del punto de extensión de idiomas (ADR-021). \\
    Turborepo~\cite{turborepo} y pnpm & Monorepo & Paquetes desacoplados con dependencias explícitas y caché de tareas (RNF-04). \\
    Vitest~\cite{vitest}, pgTAP~\cite{pgtap} y Playwright~\cite{playwright} & Pruebas & Unitarias, de base de datos (incluidas pruebas de ataque a las políticas RLS) y de extremo a extremo (RNF-06). \\
    Docker~\cite{docker-docs} & Entorno local & Ejecuta la instancia local de Supabase. \\
    Git, GitHub y GitHub Actions~\cite{github-actions} & Versiones e integración continua & Repositorio privado y CI con controles del TFG, tipos, \textit{lint}, pruebas unitarias, pgTAP y E2E. \\
    Claude Code~\cite{claudecode} & Asistente de programación & Desarrollo asistido por agentes con el arnés de la sección~\ref{sec:plan-ia}. \\
    \LaTeX{} & Memoria & Plantilla oficial de la ESI. \\
    \bottomrule
  \end{tabularx}
\end{table}

% ---------------------------------------------------------------------
\section{Planificación del proyecto}
\label{sec:plan-planificacion}

\subsection{Fases y relación con la EDT}
\label{sec:plan-fases}

La propuesta define una EDT con seis paquetes de trabajo: (1)~Análisis,
(2)~Arquitectura, (3)~Modelo de datos, (4)~Implementación, (5)~Validación y
(6)~Documentación. La tabla~\ref{tab:plan-fases} relaciona con ellos cada fase
del proyecto y su criterio de «hecho».

\begin{table}[htbp]
  \centering
  \small
  \caption{Fases del proyecto, paquetes de la EDT que cubren y criterio de «hecho».}
  \label{tab:plan-fases}
  \begin{tabularx}{\textwidth}{@{}l>{\raggedright\arraybackslash}p{3.4cm}cX@{}}
    \toprule
    \textbf{Fase} & \textbf{Nombre} & \textbf{EDT} & \textbf{Hecho cuando\ldots} \\
    \midrule
    F0 & Harness y planificación & 1 & Los controles automáticos pasan, los \textit{hooks} funcionan y existe el primer \textit{commit} \\
    F1 & Base de la aplicación web funcionando & 4 & La aplicación arranca en local, CU-01 funciona a mano y la CI está en verde \\
    F2 & Integración del CMS & 2, 3, 4 & Cada incremento F2.1--F2.9 cumple su criterio (BD, API, interfaz, explorador, usuarios y almacenamiento, auditoría, ajustes y RBAC, paneles, cierre) \\
    F3 & Interfaz en español e i18n & 4 & Interfaz en español por defecto y controles automáticos sin excepciones \\
    F4 & Comentarios didácticos en español & 4, 6 & Sin comentarios en inglés en los módulos críticos y revisión aprobada \\
    F5 & Adaptación a pymes y configurabilidad & 2, 3, 4 & CU-02 a CU-05 funcionan y la configuración está documentada \\
    F6 & Validación & 5 & CI en verde y resultados documentados \\
    F7 & Despliegue & 4, 6 & Despliegue de demostración accesible y manual probado desde cero \\
    F8 & Memoria y manuales & 1, 2, 3, 6 & Continua; se cierra tras F6 y F7 \\
    \bottomrule
  \end{tabularx}
\end{table}

Las dependencias principales son lineales (F0 $\rightarrow$ F1 $\rightarrow$
F2 $\rightarrow$ F3 $\rightarrow$ F5 $\rightarrow$ F6), con dos excepciones: F4
puede solaparse con F3 y F5, y F8 se desarrolla de forma continua, de modo que
cada fase cerrada deja un borrador de su parte de la memoria.

La fase F2 es la de mayor riesgo. Al comenzarla se evaluaron tres formas de
integrar el CMS ---aplicaciones separadas, integración híbrida e integración
total--- y se eligió la integración total en la consola de administración
(ADR-011), a pesar de ser la de mayor coste estimado, por dar lugar a un único
servicio, un único inicio de sesión y un único estilo de código. Para
controlar ese riesgo, la fase se dividió en incrementos y se identificó de
antemano el candidato a recortar si el plazo lo exigía: los paneles
configurables (RF-11, requisito deseable). Finalmente no hizo falta
recortarlo.

\subsection{Calendario real}
\label{sec:plan-calendario}

La figura~\ref{fig:plan-gantt} muestra el diagrama de Gantt de lo ejecutado
hasta la fecha de corte (20/10/2026), reconstruido a partir de las fechas de
los \textit{commits} y del registro de progreso. La
tabla~\ref{tab:plan-hitos} recoge los hitos.

\begin{figure}[htbp]
  \centering
  \resizebox{\textwidth}{!}{%
  \begin{ganttchart}[
      time slot format=isodate,
      x unit=0.55cm,
      y unit chart=0.62cm,
      y unit title=0.55cm,
      vgrid={*1{gray!25, dotted}},
      title label font=\footnotesize,
      bar label font=\footnotesize,
      group label font=\footnotesize\bfseries,
      milestone label font=\footnotesize\itshape,
      bar/.append style={fill=blue!45},
      group/.append style={fill=gray!60},
      milestone/.append style={fill=red!70},
      bar height=0.55,
      group peaks height=0.2,
    ]{2026-09-29}{2026-10-20}
    \gantttitle{sep.}{2}
    \gantttitle{octubre de 2026}{20} \\
    \gantttitlecalendar{day} \\
    \ganttbar{F0 · Harness y plan}{2026-09-29}{2026-09-29} \\
    \ganttbar{F1 · Base de la web}{2026-09-29}{2026-09-29} \\
    \ganttmilestone{Fusión de F1 en \texttt{main}}{2026-09-29} \\
    \ganttgroup{F2 · Integración del CMS}{2026-09-29}{2026-10-02} \\
    \ganttbar{F2.1--F2.4 · BD, API, interfaz, explorador}{2026-09-29}{2026-09-29} \\
    \ganttbar{F2.5--F2.7b · usuarios, archivos, auditoría, ajustes}{2026-09-30}{2026-09-30} \\
    \ganttbar{F2.7c--F2.9 · recursos, paneles, cierre}{2026-10-02}{2026-10-02} \\
    \ganttmilestone{Fusión de F2}{2026-10-02} \\
    \ganttgroup{F3 · Español y consola por áreas}{2026-09-30}{2026-10-05} \\
    \ganttbar{Identidad visual (adelantada)}{2026-09-30}{2026-09-30} \\
    \ganttbar{F3a--F3c · español, consola por áreas}{2026-10-02}{2026-10-05} \\
    \ganttmilestone{Fusión de F3 (PR \#2)}{2026-10-02} \\
    \ganttbar{Bloque académico en la landing}{2026-10-20}{2026-10-20} \\
    \ganttbar[bar/.append style={fill=green!40}]{F8 · Memoria (continua)}{2026-09-29}{2026-10-20}
  \end{ganttchart}%
  }
  \caption[Diagrama de Gantt de la ejecución real]{Diagrama de Gantt de la
    ejecución real hasta el 20/10/2026. Las barras abarcan días naturales; el
    trabajo efectivo se concentró en los días 29/09, 30/09, 02/10, 05/10 y
    20/10. Las fases F4--F7 no se habían iniciado en la fecha de corte.}
  \label{fig:plan-gantt}
\end{figure}

\begin{table}[htbp]
  \centering
  \small
  \caption{Hitos del proyecto según el historial del repositorio.}
  \label{tab:plan-hitos}
  \begin{tabularx}{\textwidth}{@{}llX@{}}
    \toprule
    \textbf{Fecha} & \textbf{\textit{Commit}} & \textbf{Hito} \\
    \midrule
    29/09/2026 & \texttt{0053dd7} & Primer \textit{commit}: directrices, plan y arnés del TFG (cierre de F0) \\
    29/09/2026 & \texttt{d9ff441} & Base de la aplicación web operativa y fusionada en \texttt{main} (cierre de F1) \\
    29/09/2026 & \texttt{48ddd1c} & Decisión de integrar el CMS en la consola (ADR-011 a ADR-014) \\
    29/09/2026 & \texttt{fd9a16a} & Primeros borradores de la memoria (planificación, análisis y diseño) \\
    02/10/2026 & \texttt{8d74ca3} & CMS integrado en la consola, fusionado en \texttt{main} (cierre de F2) \\
    02/10/2026 & \texttt{dc75f6e} & Español por defecto, fusionado (cierre de F3a--F3b) \\
    02/10/2026 & \texttt{ab1ee4c} & Interfaz solo en español (ADR-021) \\
    05/10/2026 & \texttt{6899a70} & Resumen en cada pantalla de la consola \\
    20/10/2026 & \texttt{75ff3fa} & Bloque académico y recorrido de la demostración en la landing \\
    \bottomrule
  \end{tabularx}
\end{table}

\todo{Completar el calendario con las fases F4--F8 cuando se ejecuten (fechas
de inicio y de cierre) y actualizar la figura~\ref{fig:plan-gantt}.}

\subsection{Dedicación registrada}
\label{sec:plan-dedicacion}

La dedicación se registra en el documento de progreso del proyecto con dos
medidas: las \textbf{horas del autor}, que son su dedicación real declarada y
las que cuentan para la planificación y los costes; y el \textbf{tiempo de
reloj}, estimado a partir de la hora de los \textit{commits}, que incluye
periodos en los que los agentes trabajan mientras el autor no está dedicado a
la tarea y es, por tanto, una cota superior. La
tabla~\ref{tab:plan-dedicacion} recoge los datos disponibles.

\begin{table}[htbp]
  \centering
  \small
  \caption{Dedicación registrada por jornada.}
  \label{tab:plan-dedicacion}
  \begin{tabularx}{\textwidth}{@{}lrrX@{}}
    \toprule
    \textbf{Fecha} & \textbf{Autor (h)} & \textbf{Reloj (h)} & \textbf{Fases trabajadas} \\
    \midrule
    29/09/2026 & 8,0 & $\approx$ 8,0 & F0, F1, F2.1--F2.4c, bitácora y borradores de la memoria \\
    30/09/2026 & 9,0 & --- & F2.5, F2.6, F2.6b, F2.7a, F2.7b, identidad visual (F3 adelantada) \\
    02/10/2026 & 6,0 & --- & F2.7c, F2.8, F2.9 (cierre de F2), F3a, F3b y F3c \\
    05/10/2026 & \textit{pendiente: h} & --- & Resumen en cada pestaña de la consola \\
    20/10/2026 & \textit{pendiente: h} & --- & Bloque académico de la landing y redacción de la memoria \\
    \midrule
    \textbf{Total registrado} & \textbf{23,0} & & F0--F3 completas \\
    \bottomrule
  \end{tabularx}
\end{table}

El 29/09 se desglosó además por tramos a partir de los \textit{commits}: F0
ocupó unos 45 minutos de reloj, F1 otros 45 y la integración del CMS
(F2.1--F2.4c) las seis horas y media restantes.

\todo{Registrar la dedicación de las jornadas posteriores al 02/10 y de las
fases F4--F8, e incluir el tiempo que no deja \textit{commits}: revisión,
pruebas manuales, reuniones con el tutor y redacción de la memoria.}

\subsection{Comparación entre lo estimado y lo real}
\label{sec:plan-estimado-real}

El plan del proyecto fijó desde el principio el orden de las fases, sus
entregables y sus criterios de «hecho», pero no una estimación en horas de
cada una. La única estimación cuantitativa registrada es la del análisis de
alternativas para integrar el CMS (ADR-011), hecha al comenzar la F2 en
términos de esfuerzo de desarrollo convencional. La
tabla~\ref{tab:plan-estimado-real} la compara con lo ocurrido.

\begin{table}[htbp]
  \centering
  \small
  \caption{Estimación registrada de la integración del CMS frente a la ejecución real.}
  \label{tab:plan-estimado-real}
  \begin{tabularx}{\textwidth}{@{}>{\raggedright\arraybackslash}p{4.2cm}>{\raggedright\arraybackslash}p{3cm}X@{}}
    \toprule
    \textbf{Alternativa (ADR-011)} & \textbf{Estimación} & \textbf{Real} \\
    \midrule
    (a) Aplicaciones separadas & 3--6 días & No ejecutada \\
    (c) Integración híbrida & 1--2 semanas & No ejecutada \\
    (b) Integración total (elegida) & 6--12 o más semanas & Cuatro días naturales (29/09 a 02/10) con tres jornadas de trabajo, que en total suman 23 h del autor e incluyen también F0, F1 y la mayor parte de F3 \\
    \bottomrule
  \end{tabularx}
\end{table}

La diferencia es de un orden de magnitud y requiere una lectura cuidadosa.
Primero, la estimación suponía un desarrollo convencional, mientras que la
ejecución fue asistida por agentes que trabajaban en paralelo y de forma
continuada durante cada jornada. Segundo, las horas del autor miden su
dedicación, no el trabajo total realizado por los agentes, por lo que no son
comparables uno a uno con una estimación de esfuerzo humano. Tercero, la
integración no está completa en todos sus detalles: el registro de progreso
mantiene una lista de pendientes menores de cada incremento, y la fase de
comentarios didácticos (F4), que afecta también al código del CMS, sigue
abierta. Aun con esas salvedades, el dato indica que el desarrollo asistido
desplaza el cuello de botella desde la escritura del código hacia la
verificación y la revisión, que es donde se concentraron las incidencias de la
bitácora. Este efecto se valora en el capítulo~\ref{cap:conclusiones}.

\todo{Añadir la carga prevista del TFG según el plan de estudios (créditos
ECTS y su equivalencia en horas, entre 25 y 30 h por crédito) y, si el autor
dispone de ella, la planificación inicial por fases (fechas o semanas
previstas) para completar la comparación.}

% ---------------------------------------------------------------------
\section{Organización}
\label{sec:plan-organizacion}

\subsection{Personas y roles}

El proyecto lo ejecuta una sola persona, que asume varios roles, con la
supervisión del tutor académico (tabla~\ref{tab:plan-roles}).

\begin{table}[htbp]
  \centering
  \small
  \caption{Roles del proyecto.}
  \label{tab:plan-roles}
  \begin{tabularx}{\textwidth}{@{}>{\raggedright\arraybackslash}p{3.8cm}>{\raggedright\arraybackslash}p{3.2cm}X@{}}
    \toprule
    \textbf{Rol} & \textbf{Quién} & \textbf{Responsabilidades} \\
    \midrule
    Director del proyecto y responsable del producto & Autor & Alcance, prioridades, decisiones de arquitectura (ADR) y aceptación de cada incremento \\
    Analista y arquitecto & Autor (asistido) & Requisitos, casos de uso, arquitectura, modelo de datos \\
    Desarrollador y probador & Autor y agentes de IA & Implementación, pruebas unitarias, de base de datos y E2E, revisión adversarial \\
    Redactor & Autor (asistido) & Documentación técnica, manuales y memoria \\
    Tutor académico & Juan Carlos de la Torre Macías & Supervisión, validación del alcance (punto abierto P-01) y revisión de la memoria \\
    \bottomrule
  \end{tabularx}
\end{table}

\subsection{Recursos materiales}

\begin{itemize}
  \item \textbf{Servidor de desarrollo remoto} con Linux y Docker, sobre el
        que se ejecutan la instancia local de Supabase, la aplicación y los
        agentes. Se accede a él mediante túnel SSH; sus servicios locales están
        aislados de Internet (ADR-010).
        \todo{Indicar las características del servidor (CPU, memoria,
        proveedor) y del equipo personal del autor.}
  \item \textbf{Servicios en la nube}: GitHub (repositorio privado e
        integración continua), Stripe en modo de pruebas y, para el
        despliegue, una plataforma de aplicaciones y un proyecto de Supabase
        en la nube (véase el capítulo de Despliegue del Sistema).
  \item \textbf{Software}: Node.js, pnpm, Turborepo, Supabase CLI, Docker,
        Vitest, pgTAP, Playwright y \LaTeX{}, todos de código abierto, y el
        asistente de programación con IA.
\end{itemize}

% ---------------------------------------------------------------------
\section{Costes}
\label{sec:plan-costes}

El coste del proyecto se estima como la suma del coste de personal y del de
los recursos materiales. El de personal se obtiene multiplicando las horas
dedicadas por el coste por hora de un ingeniero informático de perfil
\textit{junior} en España, incluidas las cotizaciones a cargo de la empresa. El
de los recursos materiales incluye la amortización del equipo, los servicios
contratados durante el proyecto y las licencias. La
tabla~\ref{tab:plan-costes} recoge la estructura del presupuesto.

\begin{table}[htbp]
  \centering
  \small
  \caption{Presupuesto del proyecto.}
  \label{tab:plan-costes}
  \begin{tabularx}{\textwidth}{@{}X>{\raggedleft\arraybackslash}p{3.8cm}r@{}}
    \toprule
    \textbf{Concepto} & \textbf{Base de cálculo} & \textbf{Importe (€)} \\
    \midrule
    Personal: horas del autor $\times$ coste por hora & \textit{pendiente: horas totales y tarifa} & \textit{pendiente} \\
    Equipo de desarrollo (amortización) & \textit{pendiente: valor y meses} & \textit{pendiente} \\
    Servidor de desarrollo remoto & \textit{pendiente: meses} & \textit{pendiente} \\
    Licencias de componentes de software & \textit{pendiente: importe} & \textit{pendiente} \\
    Asistente de programación con IA & \textit{pendiente: meses y plan} & \textit{pendiente} \\
    Servicios para el despliegue (plataforma y base de datos en la nube) & \textit{pendiente: plan} & \textit{pendiente} \\
    Otros (dominio, etc.) & \textit{pendiente} & \textit{pendiente} \\
    \midrule
    \textbf{Total} & & \textit{pendiente} \\
    \bottomrule
  \end{tabularx}
\end{table}

Las herramientas de desarrollo empleadas son de código abierto y no tienen
coste de licencia, y Stripe se usa en modo de pruebas, sin coste. Conviene
subrayar que, en un escenario profesional, el ahorro que justifica una base
reutilizable no se aprecia en este presupuesto sino en el de los proyectos
siguientes: cada SaaS que se construye sobre PymeKit evita el coste de
desarrollar y verificar de nuevo sus componentes comunes (véase la evaluación
de la reutilización en el anexo~\ref{cap:evaluacion}).
\todo{Completar los importes de la tabla~\ref{tab:plan-costes}, indicar la fuente de la tarifa horaria (convenio o estudio salarial) y
confirmar los planes de GitHub, Supabase y la plataforma de despliegue.}

% ---------------------------------------------------------------------
\section{Riesgos}
\label{sec:plan-riesgos}

La tabla~\ref{tab:plan-riesgos} recoge los riesgos identificados, su
mitigación y si se han materializado, con las incidencias de la bitácora que
lo documentan (identificadores B-xx).

\begin{table}[htbp]
  \centering
  \footnotesize
  \caption{Riesgos del proyecto, mitigación y materialización.}
  \label{tab:plan-riesgos}
  \begin{tabularx}{\textwidth}{@{}l>{\raggedright\arraybackslash}p{3.5cm}>{\raggedright\arraybackslash}X>{\raggedright\arraybackslash}p{2.3cm}@{}}
    \toprule
    \textbf{ID} & \textbf{Riesgo} & \textbf{Mitigación} & \textbf{Materializado} \\
    \midrule
    R-01 & Desbordamiento del plazo por la integración del CMS &
      Incrementos funcionales F2.1--F2.9; RF-11 identificado como primer recorte (ADR-011) & No \\
    R-02 & Fallos de seguridad en el control de acceso &
      Auditoría adversarial de RLS con pruebas pgTAP de ataque y refutación independiente; pruebas de regresión (ADR-015) & Sí: B-05--B-09, B-19--B-25, B-27, B-28, B-42, B-47, B-50 \\
    R-03 & Exposición del entorno de desarrollo (servidor con IP pública) &
      Filtrado en la cadena \texttt{DOCKER-USER}, acceso por túnel SSH y revisión de cada servicio nuevo (ADR-010) & Sí: B-04 \\
    R-04 & Errores introducidos por cambios masivos o por los agentes &
      Controles automáticos que cubren las variantes del error, revisión adversarial y prueba manual del autor & Sí: B-01, B-14, B-16, B-45 \\
    R-05 & Deriva entre el esquema declarativo y las migraciones &
      Diferencia de esquemas en la verificación de todo cambio de BD; propiedades de seguridad comprobadas con pgTAP & Sí: B-10, B-35 \\
    R-06 & Pruebas que no prueban lo que dicen o que quedan desactualizadas &
      Leer las aserciones de las pruebas existentes; suite completa antes de fusionar & Sí: B-11, B-51, B-54 \\
    R-07 & Interrupción de las herramientas o de la conexión durante una tarea &
      Cerrar cada incremento con su \textit{commit}; tareas de seguridad formuladas como verificación defensiva y por etapas & Sí: B-12, B-15 \\
    R-08 & Falsos positivos de las herramientas (secretos, pruebas inestables) &
      Documentar las credenciales públicas de desarrollo; E2E contra la \textit{build} de test & Sí: B-03, B-13 \\
    R-09 & Coste excesivo de la verificación asistida por agentes &
      Modo de ahorro: un agente por paso y suite completa al cerrar la fase & Sí: B-48 \\
    R-10 & Incumplimiento de las licencias de las dependencias &
      Repositorio privado y revisión de las condiciones de licencia (ADR-001, ADR-006) & No \\
    R-11 & Validación de la reutilización sin un dominio de negocio concreto &
      Medir el arranque de un SaaS nuevo a partir de la plataforma (ADR-005, P-01) & Pendiente \\
    \bottomrule
  \end{tabularx}
\end{table}

\todo{Asignar probabilidad e impacto (bajo, medio o alto) a cada riesgo según
el criterio del autor y, si se desea, una matriz de riesgos.}

Tres riesgos materializados merecen comentario por su enseñanza para la
planificación:

\begin{itemize}
  \item \textbf{B-04, exposición del entorno.} Los servicios locales de la
        base de datos publicaban sus puertos en todas las interfaces y el
        motor de contenedores insertaba reglas que se saltaban el cortafuegos,
        de modo que la base de datos de desarrollo, con credenciales conocidas
        y solo datos de prueba, fue accesible desde Internet durante
        aproximadamente una hora. Se corrigió con una regla persistente, se
        recrearon los volúmenes y se pasó a trabajar mediante un túnel SSH
        (ADR-010).
  \item \textbf{Fallos de control de acceso.} Las auditorías de la base de
        datos del CMS encontraron brechas reales, entre ellas una
        escalada de privilegios en el RBAC (B-05), comprobaciones que fallaban
        en abierto (B-06, B-23), una inyección SQL (B-24) y un escape de las
        denegaciones explícitas mediante delegación (B-47). El riesgo estaba
        previsto, y la mitigación ---auditoría obligatoria en todo cambio de
        base de datos--- es la que permitió detectarlos.
  \item \textbf{B-15, corte de conexión con un agente trabajando.} El trabajo
        quedó a medias y mezclado con una corrección de seguridad pendiente. La
        lección pasó al proceso: cada incremento se cierra con su propio
        \textit{commit}.
\end{itemize}

% ---------------------------------------------------------------------
\section{Gestión de la configuración}
\label{sec:plan-configuracion}

\begin{itemize}
  \item \textbf{Control de versiones.} Git, con un repositorio privado en
        GitHub (ADR-001). La
        rama principal es \texttt{main}; se trabaja en ramas por fase
        (\texttt{fase-1/base}, \texttt{fase-2/cms}, \texttt{fase-3/espanol}) que
        se fusionan en \texttt{main} al cerrar la fase, mediante
        \textit{pull request} cuando se requiere ejecutar la suite E2E en la
        CI. Los mensajes de \textit{commit} están en español y siguen
        \textit{Conventional Commits}~\cite{conventionalcommits} con el módulo
        como ámbito, por ejemplo \texttt{feat(cms): auditoría y búsqueda
        global (F2.6)}.
  \item \textbf{Base de datos como código.} El esquema se mantiene de forma
        declarativa en ficheros numerados, y las migraciones se generan a
        partir de ellos y se versionan en el repositorio. Nunca se modifica una
        migración ya aplicada: cada cambio es una migración nueva y no
        destructiva. Las pruebas pgTAP viven junto al esquema.
  \item \textbf{Configuración y secretos.} La configuración reutilizable está
        en ficheros \texttt{*.config.ts} validados con Zod y en variables de
        entorno públicas documentadas; los secretos solo se guardan en
        ficheros locales excluidos del repositorio o en las variables de la
        plataforma de despliegue.
  \item \textbf{Dependencias.} Un catálogo único de versiones para todo el
        monorepo, con el fichero de bloqueo versionado.
  \item \textbf{Entregas.} Cada fusión de una fase en \texttt{main} constituye
        una versión estable del producto (tabla~\ref{tab:plan-hitos}). El
        despliegue de demostración se construye desde \texttt{main}.
        \todo{Decidir si se etiquetan versiones (\texttt{git tag}) para la
        entrega final del TFG.}
\end{itemize}

% ---------------------------------------------------------------------
\section{Aseguramiento de calidad}
\label{sec:plan-calidad}

\subsection{Estándares y normas}

\begin{itemize}
  \item \textbf{Código}: TypeScript estricto, sin \texttt{any} sin
        justificar; \textit{lint} y formato automáticos; convenciones del
        repositorio para \textit{server functions}, formularios, clientes de
        base de datos y separación entre módulos de cliente y de servidor.
  \item \textbf{Comentarios}: guía propia de comentarios didácticos en
        español, con cabecera por fichero, documentación de las funciones
        exportadas y justificación obligatoria de toda decisión de seguridad.
  \item \textbf{Base de datos}: RLS obligatoria en toda tabla nueva,
        privilegios mínimos por verbo y por columna, y ninguna función con
        privilegios elevados sin controles explícitos.
  \item \textbf{Seguridad}: el estándar OWASP ASVS~\cite{owasp-asvs} sirve de
        referencia para los controles de autenticación, sesión y control de
        acceso.
  \item \textbf{Calidad del producto}: los requisitos no funcionales se
        alinean con las características del modelo ISO/IEC
        25010~\cite{iso25010}: seguridad (RNF-02, RNF-03), mantenibilidad y
        modularidad (RNF-04, RNF-05), portabilidad e instalabilidad (RNF-07),
        usabilidad (RNF-08) y eficiencia (RNF-09).
  \item \textbf{Trazabilidad}: una matriz relaciona cada requisito con el
        código que lo implementa, las pruebas que lo verifican y la sección de
        la memoria que lo describe. Los puntos clave del código llevan
        etiquetas \texttt{[TFG]} que un script extrae a esa matriz.
\end{itemize}

\subsection{Verificación y validación}

Todo cambio pasa por la misma cadena de comprobaciones, declarada obligatoria
en las directrices del repositorio:

\begin{enumerate}
  \item comprobación de tipos, \textit{lint} y formato;
  \item controles automáticos del repositorio, tanto sobre el contenido
        como sobre las rutas de los ficheros;
  \item pruebas afectadas: unitarias, de base de datos (pgTAP) y E2E si se
        toca un flujo de usuario;
  \item revisión adversarial del cambio en dos pistas: calidad y
        arquitectura, y búsqueda de errores y fallos de seguridad;
  \item si el cambio afecta a la base de datos, auditoría adversarial de las
        políticas RLS en cuatro etapas: revisión estática, pruebas de ataque
        ejecutables entre clientes, refutación independiente de los hallazgos
        y corrección con pruebas de regresión;
  \item revisión de los comentarios del código; y
  \item actualización del progreso, la trazabilidad y, si procede, las
        decisiones y la bitácora.
\end{enumerate}

Las comprobaciones automatizables se repiten en la integración continua, en un
\textit{workflow} con tres trabajos: (1)~controles propios del TFG
(controles automáticos del repositorio, bloqueantes; informe de comentarios en inglés, informativo hasta
cerrar la F4); (2)~tipos, \textit{lint}, formato y pruebas unitarias; y
(3)~base de datos local, pruebas pgTAP y E2E con Playwright contra la
\textit{build} de test, que se ejecuta en las \textit{pull requests} hacia
\texttt{main} y bajo demanda.

\subsection{Criterios de aceptación y acciones correctoras}

Un incremento se acepta cuando cumple su criterio de «hecho», la cadena
anterior está en verde y el autor lo ha revisado. Se rechaza ---y se corrige
antes de continuar--- si una auditoría encuentra una brecha confirmada, si
alguna prueba deja de pasar o si los controles automáticos del repositorio
emiten un aviso. Cada incidencia se registra en la bitácora con su lección, y
cuando la lección es generalizable se convierte en una acción preventiva: una
regla nueva en las directrices, una variante más en un control automático
(B-01) o una prueba de regresión. A la fecha de corte la bitácora
recoge 55 incidencias, de las que 29 son de seguridad; todas las de gravedad alta están corregidas.
